\documentclass[conference]{IEEEtran}
\IEEEoverridecommandlockouts
\usepackage{cite}
\usepackage{amsmath,amssymb,amsfonts}
\usepackage{graphicx}
\usepackage{textcomp}
\usepackage{xcolor}
\usepackage{booktabs}

% ---------------------------------------------------------------------------
% Drafting aids. Set \drafttrue to show outline notes in the PDF,
% \draftfalse to hide all of them before submission.
% ---------------------------------------------------------------------------
\newif\ifdraft
\drafttrue
\newcommand{\note}[1]{\ifdraft{\color{blue}\footnotesize [#1]}\fi}
\newcommand{\todo}[1]{\ifdraft{\color{red}\footnotesize [TODO: #1]}\fi}
% Placeholder box for a figure that has not been added yet.
\newcommand{\figplaceholder}[1]{%
  \fbox{\parbox[c][0.45\columnwidth][c]{0.9\columnwidth}{\centering\footnotesize #1}}}

% ---------------------------------------------------------------------------
% Notation (keep in sync with claude/handover_paper_kalman_afc_fan.md, Sec. 0)
% ---------------------------------------------------------------------------
\newcommand{\Ts}{T_s}
\newcommand{\thc}{\theta_c}              % AFC state p - jq
\newcommand{\thch}{\hat{\theta}_c}       % shadow AFC state
\newcommand{\etad}{\eta_d}               % envelope AR(1) coefficient
\newcommand{\Gp}{G'}                     % loop gain at omega_0 incl. direct term
\newcommand{\E}{\mathrm{E}}
\newcommand{\re}{\operatorname{Re}}

\begin{document}

% Tentative title -- to be fixed after the draft is complete.
\title{Shadow Adaptive Feedforward Cancellation Based on Kalman Estimation for
Fan-Induced Vibration After Track Seek in Hard Disk Drives}

\author{\IEEEauthorblockN{Shota Yabui}
\IEEEauthorblockA{\textit{Tokyo City University}\\
Tokyo, Japan}
\and
\IEEEauthorblockN{Takenori Atsumi}
\IEEEauthorblockA{\textit{Chiba Institute of Technology}\\
Chiba, Japan}
}

\maketitle

\begin{abstract}
\todo{Write last, about 150--200 words. No symbols or math.}
\note{Flow: (1) fan-induced vibration in HDDs is a random narrow-band
disturbance; (2) disturbance-rejection filters are frozen during track seek, so
the disturbance estimate is lost and the write-ready time is delayed; (3) the
learned coefficients of AFC are interpreted as a Kalman prediction of the
disturbance envelope; (4) this gives the time update during seek and a shadow
AFC that keeps estimating without driving the plant; (5) benchmark simulation
with 1000 disturbance sequences shows a shorter write-ready time with the
disturbance-free settling unchanged.}
\end{abstract}

\begin{IEEEkeywords}
hard disk drive, adaptive feedforward cancellation, Kalman filter,
narrow-band disturbance, track seek, dual-stage actuator
\end{IEEEkeywords}

% ===========================================================================
\section{Introduction}
% Target length: 0.75 page
% ===========================================================================

\note{P1 Background. Head positioning in HDDs with a dual-stage actuator
(VCM + PZT). Cooling-fan vibration in storage enclosures is a dominant
disturbance in the kHz range.}

\note{P2 Existing approach. Resonant filters / AFC shaped by loop shaping
suppress narrow-band disturbances during track following
\cite{bodson1994,messner1995,atsumi2026amc,yabui_thesis}. Parameters are chosen
from the sensitivity function, not from the statistics of the disturbance.}

\note{P3 Problem. The fan vibration is random narrow-band, not periodic: its
envelope drifts with a correlation time $1/(\zeta_d\omega_0)$. During a seek
the filter input is gated (or the state is reset), so the estimate decays and
must be re-learned after mode switching. This delays the write-ready time.}

\note{P4 Idea. The filters themselves stay loop-shaped. We add an
estimation-theoretic view: the AFC state is an estimate of the complex envelope
of the disturbance, and a steady-state Kalman predictor tells us (i) how the
estimate should evolve while no measurement is used and (ii) how to keep using
the position error during seek without disturbing the seek response.}

\note{P5 Contributions (three items):
(1) interpretation of the AFC learned coefficients as an envelope estimate, and
its relation to the MMSE (Kalman) predictor: a loop-shaped AFC gives a shrunk
and rotated estimate with a short memory;
(2) shadow AFC that separates estimation from actuation during seek, with a
seek-time gain that treats model error as measurement noise;
(3) paired evaluation on the IEEJ HDD benchmark with 1000 disturbance
sequences.}

% ===========================================================================
\section{Problem Setting}
% Target length: 0.9 page. Fig. 1 (disturbance PSD), Fig. 2 (block diagram)
% ===========================================================================

\subsection{Benchmark and Control System}
\note{IEEJ HDD benchmark, Track-Seek Control ver.~1.0.3 \cite{ieej_benchmark}.
Sampling period $\Ts = 19.84~\mu$s. Three plant cases LT/RT/HT. Mode-switching
feedback: seek controller during the trajectory, switched to the following
controller with disturbance-rejection filters 0.3~ms after the trajectory ends.
Feedforward: polynomial trajectory FF.}

\todo{State the number of disturbance-rejection filters in the following
controller used for the results (handover says 13 for FB v2; the MATLAB design
tool gives 12). All of them are loop-shaped.}

\begin{figure}[tb]
\centering
\figplaceholder{Block diagram of the mode-switching dual-stage control
system with AFC}
\caption{Control system. \todo{draw}}
\label{fig:block}
\end{figure}

\subsection{Fan-Induced Vibration}
\note{Sum of 13 second-order resonances (2.2--9.2~kHz), each driven by white
noise: random narrow-band disturbance. Focus on 6.66~kHz ($\zeta_d = 0.008$,
correlation time 3.0~ms): largest in disturbance PSD, disturbance $\sigma$
(1.88~nm) and position error without filters (1.52~nm).}

\begin{figure}[tb]
\centering
\figplaceholder{PSD of the fan-induced disturbance (13 modes) and PES\\
source: designs/20261106\_fan\_psd/fan\_psd.png}
\caption{Power spectral density of the fan-induced vibration. \todo{add figure}}
\label{fig:psd}
\end{figure}

\subsection{AFC and Its Handling During Seek}
The AFC for the angular frequency $\omega_0$ ($\Omega = \omega_0 \Ts$) generates
\begin{equation}
u(k) = \re\{\thc(k) e^{j\Omega k}\}, \qquad \thc = p - jq ,
\label{eq:afc_out}
\end{equation}
and updates its state from the input $x(k)$ by
\begin{equation}
\thc(k+1) = \eta\,\thc(k) + \lambda\, x(k)\, e^{-j(\Omega k + \theta)} ,
\label{eq:afc_update}
\end{equation}
where $\lambda$, $\theta$, and $\eta$ are the learning gain, the learning phase,
and the forgetting factor.
\note{Equivalent to a resonant filter with damping $\zeta_a$
\cite{bodson1994,yabui_thesis}. Add direct term $n_0$ if space allows.}

\note{Conventional handling during seek: \emph{gate} (input set to zero, so
$\thc \leftarrow \eta\,\thc$) or \emph{reset} ($\thc = 0$). Evaluation index:
write-ready time $t_\mathrm{ready}$ = first time the PES stays within
$\pm 0.1$~track for 0.5~ms, measured from seek start. Disturbance-free settling
time $t_\mathrm{hold}$ and $\max|u_{cp}| < 50$~nm must not be degraded.}

\note{Key observation: the reference is the predicted position of the nominal
model, so the PES during seek contains only disturbance and model error. The
disturbance is observable during seek.}

% ===========================================================================
\section{Kalman Interpretation of AFC}
% Target length: 1.25 pages. Table I (6.66 kHz), Fig. 3 (envelope estimate)
% ===========================================================================

\subsection{Envelope Model of the Disturbance}
The disturbance at the AFC input is written with a complex envelope $D(k)$ as
\begin{equation}
d(k) = \re\{D(k) e^{j\Omega k}\}, \qquad \E|D|^2 = 2\sigma_d^2 ,
\end{equation}
and the envelope of a lightly damped mode is approximated by
\begin{equation}
D(k+1) = \etad D(k) + v(k), \qquad \etad = e^{-\zeta_d \omega_0 \Ts},
\label{eq:ar1}
\end{equation}
with $Q = \E|v|^2 = 2\sigma_d^2(1-\etad^2)$.

\subsection{Demodulated Input as Innovation}
Let $\vartheta = -\Gp(\omega_0)\thc$ be the disturbance component cancelled by
the AFC, where $\Gp$ is the loop from the AFC output to its input. Demodulating
the input gives
\begin{equation}
y(k) = 2x(k)e^{-j\Omega k} \simeq D(k) - \vartheta(k) + w(k),
\label{eq:innov}
\end{equation}
with $R = \E|w|^2 = 4\sigma_n^2$, where $\sigma_n^2$ is the equivalent variance
of the other disturbance components at $\omega_0$.
\note{$y$ is ``measurement minus prediction'', i.e.\ the innovation. Image
component at $2\Omega$ neglected because learning is slow.}

\subsection{Steady-State Kalman Predictor}
The steady-state Kalman predictor \cite{kalman1960,anderson1979} for
\eqref{eq:ar1} and \eqref{eq:innov} is
\begin{equation}
\vartheta(k+1) = \etad\,\vartheta(k) + \beta\, y(k), \qquad
\beta = \frac{\etad P}{P+R},
\label{eq:kalman}
\end{equation}
\begin{equation}
P = \frac{-b + \sqrt{b^2 + 4QR}}{2}, \qquad b = R(1-\etad^2) - Q .
\end{equation}
Comparing \eqref{eq:kalman} with \eqref{eq:afc_update} through
$\thc = -\vartheta/\Gp$ gives
\begin{equation}
\eta = \etad, \qquad
\lambda = \frac{2\beta}{|\Gp(\omega_0)|}, \qquad
\theta = \arg \Gp(\omega_0) + \pi .
\label{eq:mapping}
\end{equation}
\note{Meaning: $\eta$ = memory length of the disturbance, $\lambda$ =
confidence in the measurement, $\theta$ = loop phase compensation.
Wording: \eqref{eq:mapping} is used as a reference for interpretation; the
following controller in this paper remains loop-shaped. Do not claim that
tuning is unnecessary (per-mode design overestimates $\beta$ for wide modes
such as 5.85~kHz).}

\subsection{What a Loop-Shaped AFC Estimates}
\note{Steady-state ratio of the estimate to the true envelope:
$\kappa = -v/(1-\eta-v)$, $v = (\lambda/2)e^{-j\theta}\Gp$.
Loop-shaped filter at 6.66~kHz: $\kappa = 0.70\angle{-10^\circ}$
(simulation: regression coefficient $0.612\angle{-11.8^\circ}$).
Kalman parameters: $\kappa = \beta/(1-\etad+\beta)$ is real (no rotation;
simulation $0.748\angle{+1.5^\circ}$, optimum $0.807\angle 0^\circ$).
So the loop-shaped AFC holds a shrunk, rotated, short-memory estimate.}

\begin{table}[tb]
\caption{AFC Parameters at 6.66 kHz}
\label{tab:param}
\centering
\begin{tabular}{lcc}
\toprule
 & Loop-shaped & Kalman \eqref{eq:mapping} \\
\midrule
Forgetting factor $\eta$ & 0.9691 & 0.99338 \\
Memory length [ms] & 0.65 & 3.0 \\
Learning strength $\beta$ & 0.0638 & 0.0546 \\
Learning gain $\lambda$ & 0.18528 & 0.15871 \\
Learning phase $\theta$ [deg] & $-25.1$ & $-59.9$ \\
\bottomrule
\end{tabular}
\end{table}

\begin{figure}[tb]
\centering
\figplaceholder{Estimated envelope $-\Gp\thc$ vs.\ true envelope
(magnitude, phase error)\\
source: designs/20261112\_pq\_envelope/pq\_envelope.png}
\caption{AFC state and the envelope of the 6.66~kHz disturbance.
\todo{add figure}}
\label{fig:envelope}
\end{figure}

% ===========================================================================
\section{AFC During Track Seek}
% Target length: 1.0 page. Fig. 4 (shadow AFC structure)
% ===========================================================================

\subsection{Time Update While the Input Is Gated}
\note{Without measurements the optimal prediction is the time update
$\thc \leftarrow \etad\,\thc$ (\emph{predict}); normalized error grows as
$e(N) = \etad^{2N}e_0 + (1-\etad^{2N})$. Gating a loop-shaped AFC instead
applies $\eta^{N}$: at 6.66~kHz the state shrinks to 0.27 after 0.84~ms,
whereas the correlation of the disturbance supports $\etad^{N} = 0.76$.
Gate is therefore close to reset. If $\eta = \etad$, gate equals predict.}

\subsection{Shadow AFC: Separating Estimation From Actuation}
\note{Motivation in one sentence: closing the AFC loop during seek uses the
measurement but interacts with residual vibration and model error
(disturbance-free settling at span 256 degraded from 0.665 to 1.364~ms), so
estimation and actuation are separated.}

During seek the plant is driven by the predicted value $\thc$, while a shadow
state $\thch$ keeps learning from the corrected input
\begin{equation}
x_\mathrm{eff}(k) = x(k) + \re\{G_s(\thch(k) - \thc(k))e^{j\Omega k}\},
\label{eq:xeff}
\end{equation}
\begin{equation}
\thch(k+1) = \eta\,\thch(k)
  + \gamma_s \lambda_s\, x_\mathrm{eff}(k)\, e^{-j(\Omega k + \theta_s)},
\label{eq:shadow}
\end{equation}
where $G_s$ is the loop with the seek controller, and
$\lambda_s = \lambda\,|G_f/G_s|$, $\theta_s = \theta - \arg(G_f/G_s)$ convert
the parameters designed for the following loop $G_f$. At mode switching the
estimate is handed over, $\thc \leftarrow \thch$.
\todo{Confirm which forgetting factor is used in \eqref{eq:shadow} and for the
output during seek ($\eta$ of the filter or $\etad$).}
\note{Since the output before switching equals that of predict, $t_\mathrm{hold}$
and $\max|u_{cp}|$ are unchanged by construction.}

\begin{figure}[tb]
\centering
\figplaceholder{Structure of the shadow AFC (output path, shadow
estimator, hand-over at mode switching)}
\caption{Shadow AFC. \todo{draw}}
\label{fig:shadow}
\end{figure}

\subsection{Seek-Time Gain}
\note{Model error during seek is treated as measurement noise. With}
\begin{equation}
Z = \sum_{k=0}^{N_s-1} \eta^{N_s-1-k}\, x(k)\, e^{-j\Omega k},
\end{equation}
\begin{equation}
\nu = \frac{\max_{U} |Z_\mathrm{det}|^2}{\E|Z_\mathrm{dist}|^2}, \qquad
\gamma_s = \frac{1}{1+\nu},
\label{eq:gamma}
\end{equation}
\note{$Z_\mathrm{det}$: disturbance-free seek over a set $U$ of 12 perturbed
plants; $Z_\mathrm{dist}$: nominal plant with disturbance only (200 sequences).
Example values of $\gamma_s$ (span 64): 0.90 at 6.66~kHz, 0.48 at 5.85~kHz,
0.06 at 3.3~kHz. \todo{these example values are for MPC FF; recompute for
polynomial FF or state the condition.}}

% ===========================================================================
\section{Simulation}
% Target length: 1.5 pages. Table II (t_ready), Fig. 5 (CDF)
% ===========================================================================

\subsection{Conditions}
\note{Benchmark plant, disturbance (RV + 13 fan modes + RRO) and evaluation
unchanged. 1000 disturbance sequences; $\gamma_s$ computed from a separate set
of 200 sequences. The same sequences are used for all methods (paired
comparison). Methods: Conventional (gate), Reset, Predict, Shadow AFC
(proposed).}

\subsection{Write-Ready Time}
\begin{table}[tb]
\caption{90th Percentile of Write-Ready Time [ms] (Polynomial FF, 1000
Sequences): Conventional $\rightarrow$ Proposed}
\label{tab:tready}
\centering
\begin{tabular}{cccc}
\toprule
Span & LT & RT & HT \\
\midrule
4    & 0.692 $\rightarrow$ 0.657 & 0.577 $\rightarrow$ 0.520 & 0.591 $\rightarrow$ 0.569 \\
64   & 1.018 $\rightarrow$ 1.004 & 0.867 $\rightarrow$ 0.836 & 1.097 $\rightarrow$ 0.954 \\
256  & 1.144 $\rightarrow$ 1.099 & 1.024 $\rightarrow$ 0.965 & 1.174 $\rightarrow$ 1.065 \\
1024 & 1.525 $\rightarrow$ 1.476 & 1.470 $\rightarrow$ 1.444 & 1.700 $\rightarrow$ 1.569 \\
\bottomrule
\end{tabular}
\end{table}

\note{Proposed is equal or shorter in all 12 conditions. 99th percentile:
span 1024 HT $2.204 \rightarrow 1.858$~ms, span 4 RT $1.006 \rightarrow
0.779$~ms. Paired comparison: difference below 0.05~ms in 65--85\% of the
sequences; when it differs the proposed method wins two to three times as
often. Reset is slower than Conventional in all conditions (e.g.\ span 64 HT:
1.182~ms).}

\todo{Results for Predict and the excess-error index with polynomial FF are
not in the handover (only MPC FF). Needed to show the step
gate $\rightarrow$ predict $\rightarrow$ shadow.}

\begin{figure}[tb]
\centering
\figplaceholder{CDF of $t_\mathrm{ready}$ and paired difference
(span 64, polynomial FF)\\
source: designs/20261108\_ready\_dist\_poly/ready\_dist\_span64.png}
\caption{Distribution of the write-ready time. \todo{add figure}}
\label{fig:cdf}
\end{figure}

\subsection{Disturbance-Free Seek}
\note{$t_\mathrm{hold}$ and $\max|u_{cp}|$ coincide with those of predict for
9 spans $\times$ 13 plants (difference 0~$\mu$s, $u_{cp}$ difference
$\le 0.18$~nm).}

% ===========================================================================
\section{Discussion and Conclusion}
% Target length: 0.4 page
% ===========================================================================

\note{Limitations to state plainly:
(1) simulation only, experiments are future work;
(2) the estimation error in the HDD loop is about twice the theoretical value
(0.38 vs.\ 0.19): colored noise from the neighbouring 5.85~kHz mode,
in-band variation of $G$, interaction with the other filters;
(3) the gain is smaller with the MPC-based FF, reason not yet identified;
(4) designing all filters by \eqref{eq:mapping} independently does not beat
loop shaping in following accuracy; joint design is future work.}

\note{Conclusion: summarize the interpretation, the shadow AFC, and the
quantitative result in three sentences.}

\bibliographystyle{IEEEtran}
\bibliography{refs}

\end{document}
